\documentclass[11pt]{article}
\usepackage[margin=1in]{geometry}
\usepackage{booktabs}
\usepackage{hyperref}
\usepackage{amsmath}

\title{Tau Scaling as an Evidence-Gated Public Claim System\\
\large Source Population, Primary-Validation Gaps, and Non-Claim Locks}
\author{James Paul Jackson}
\date{TAU-SCALING-SA v0.9.3}

\begin{document}
\maketitle

\begin{abstract}
Public semiconductor narratives often combine roadmap projections, media interpretation, methodology framing, reported metrics, and implied validation. This manuscript presents a repository-based method for converting public Tau Scaling and LogicFolding claims into an evidence-gated claim system. In the current repository state, eight public Tau claims are source-populated from public reporting, zero are primary-source confirmed, and zero are independently confirmed. The contribution is a reproducible evidence-governance artifact for separating public narrative from validation status. It is not a claim of silicon validation, product validation, benchmark superiority, process-node equivalence, or a universal Tau Scaling law.
\end{abstract}

\section{Contribution}
This work contributes a reproducible repository method for public semiconductor claim governance, a source-provenance workflow, a validation-gap matrix, manuscript-ready evidence tables, and explicit non-claim locks.

\section{Research Question}
Can public Tau Scaling claims be represented as a reproducible evidence-governance system that separates source population from primary validation and claim promotion?

\section{Method}
The method converts public claims into a staged evidence pipeline: claim ledger, source provenance map, source intake card, source population manifest, primary-source validation gap review, public research milestone package, manuscript draft, evidence pack, and submission package.

\section{Results}
\begin{center}
\begin{tabular}{lr}
\toprule
Metric & Value\\
\midrule
Public claims analyzed & 8\\
Source-populated claims & 8\\
Primary-source confirmed claims & 0\\
Independent-source confirmed claims & 0\\
High validation-gap claims & 8\\
Release findings & 0\\
\bottomrule
\end{tabular}
\end{center}

\section{Interpretation}
The principal result is not validation of Tau Scaling. The principal result is the reproducible separation of public claim visibility, source population, primary-source confirmation, independent validation, and claim promotion.

\section{Limitations}
This manuscript does not validate semiconductor performance, Huawei product performance, manufacturing capability, process-node equivalence, benchmark superiority, or a universal Tau Scaling law.

\section{Data and Code Availability}
All artifacts are repository-local under docs, reports, releases, sources, visuals, and scripts.

\section{Ethics and Non-Claim Statement}
This manuscript concerns evidence governance. It should not be read as investment advice, industrial validation, product endorsement, national capability assessment, or proof of semiconductor performance.

\section{Conclusion}
The current repository state supports a narrow, publishable claim: Tau Scaling public claims can be transformed into a reproducible evidence-governance artifact with source-population and validation-gap surfaces.

\end{document}
